\subsection{作逆紧作用的局部紧群}

命题 7. 设 G 是连续作用于豪斯多夫空间 X 的局部紧群。则 G 逆紧作用于 X 当且仅当：对 X 的每一对点 x, y，存在 x 的邻域 \(V_{x}\) 与 y 的邻域 \(V_{y}\)，使得集合 K——由 \(s \in G\) 中满足 \(s \cdot V_{x} \cap V_{y} \neq \emptyset\) 的所有 s 组成——在 G 中相对紧。

设 F 是把无穷远点 \(\omega\) 添加到 G 上所得的紧空间，并设 \(\Gamma\) 是 \(\rho: (s, x) \to s.x\) 的图，把它看作 \(F \times X \times X\) 的一个子集。我们来证明，如果 \(pr_{23}\) 在 \(\Gamma\) 上的限制是逆紧的，那么 \(\Gamma\) 在 \(F \times X \times X\) 中是闭的。事实上，这一假设蕴含映射 \(u: (t, s, x, y) \to (t, x, y)\)（\(F \times \Gamma\) 到 \(F \times X \times X\) 中的映射）是闭映射。若 \(\Gamma'\) 是由点 \((s, s)\) 组成的集合，这些点位于 \(F \times G\) 中且 \(s \in G\)，则 \(\Gamma'\) 在 \(F \times G\) 中闭，因为它是典范单射 \(G \to F\) 的图（第一章，§8，第 1 节，命题 2 的推论 2）；于是交 \((\Gamma' \times X \times X) \cap (F \times \Gamma)\) 在 \(F \times \Gamma\) 中闭，并且立即可以看出它在 u 下的像正是把 \(\Gamma\) 看作 \(F \times X \times X\) 的子集时的那个集合；从而 \(\Gamma\) 在 \(F \times X \times X\) 中闭。另一方面，我们有 \((\{\omega\} \times X \times X) \cap \Gamma = \emptyset\)。于是由 F 的定义，对每一点 \((x, y) \in X \times X\)，存在 \((x, y)\) 在 \(X \times X\) 中的一个邻域 W 和 G 的一个紧子集 K，使得

\[
((\mathbf {G} - \mathbf {K}) \times \mathbf {W}) \cap \Gamma ;
\]

为空；由于可取 W 为形如 \(V_{x} \times V_{y}\) 的邻域，其中 \(V_{x}\) 与 \(V_{y}\) 分别是 X 中 x 与 y 的邻域，断言“((G - K) × W) ∩ Γ = ∅”就成为“若 s∉K，则 s.Vx∩Vy=∅”。这样我们就证明了命题所述条件的必要性。反之，设此条件成立；设 A 是由超滤子 \(\mathfrak{F}\) 滤化的集合，并设 \(\alpha \to (s_{\alpha}, x_{\alpha})\) 是 A 到 \(G \times X\) 中的映射，使得 \(\lim_{\mathfrak{F}} x_{\alpha} = x\) 且 \(\lim_{\mathfrak{F}} s_{\alpha}\cdot x_{\alpha} = y\)。设 K、Vx 与 Vy 满足命题的条件。由假设，存在集合 \(M \in \mathfrak{F}\)，使得若 \(\alpha \in M\)，则 \(x_{\alpha} \in V_{x}\) 且 \(s_{\alpha}\cdot x_{\alpha} \in V_{y}\)，从而 \(s_{\alpha} \in K\)。这表明 \(\alpha \to s_{\alpha}\) 关于 \(\mathfrak{F}\) 收敛，证毕。

若 G 是紧的，命题 7 的条件平凡地成立；这样我们重新得到命题 2a)。

命题 7 特别表明，连续作用于豪斯多夫空间 X 的离散群 G 逆紧作用于 X 当且仅当：对 X 的每一对点 \((x, y)\)，存在 x 的邻域 \(V_{x}\) 与 y 的邻域 \(V_{y}\)，使得由 \(s \in G\) 中满足 \(s \cdot V_{x} \cap V_{y} \neq \emptyset\) 的点组成的集合是有限的。

命题 8. 设 G 是逆紧作用于豪斯多夫空间 X 的离散群。设 x 是 X 的一点，\(K_{x}\) 是 x 的稳定子群。那么：

\begin{enumerate}
\def\labelenumi{\alph{enumi})}
\item
  子群 \(K_{x}\) 是有限的，并且存在 \(U \subset X\) 的开集，它包含 x，在 \(K_{x}\) 下稳定，并且在它上面由 G 所定义的关系诱导的等价关系就是由 \(K_{x}\) 定义的等价关系。
\item
  典范映射 \(\mathbf{U} / \mathbf{K}_x\to \mathbf{X} / \mathbf{G}\) 是 \(\mathbf{U} / \mathbf{K}_x\) 到 \(x\) 的类在 \(\mathbf{X} / \mathbf{G}\) 中的一个开邻域上的同胚。
\end{enumerate}

由命题 7，\(\mathbf{K}_x\) 是有限的。为构造满足所需条件的开集 \(\mathbf{U}\)，首先注意，由命题 7 存在开集 \(\mathrm{U}_0\)（包含 \(x\)），使得集合 \(\mathbf{K}\)——由 \(s \in \mathbf{G}\) 中满足 \(s. \mathrm{U}_0 \cap \mathrm{U}_0 \neq \emptyset\) 的那些 s 组成——是有限的。显然 \(\mathbf{K}_x \subset \mathbf{K}\)；设 \(s_1, \ldots, s_n\) 是 \(\mathbf{K} - \mathbf{K}_x\) 的元素。若令 \(x_i = s_i.x\)（ \(1 \leqslant i \leqslant n\)），则 \(x_i \neq x\) 对每个 \(i\) 成立；由于 \(\mathbf{X}\) 是豪斯多夫的，对每个指标 \(i\) 存在开邻域 \(\mathrm{V}_i\)（包含 \(x\)）与开邻域 \(\mathrm{V}_i'\)（包含 \(s_i.x\)），使得 \(\mathrm{V}_i \cap \mathrm{V}_i' = \emptyset\)。令 \(\mathrm{U}_i = \mathrm{V}_i \cap s_i^{-1}. \mathrm{V}_i'\)；则 \(\mathrm{U}_i\) 显然是开的且包含 \(x\)，并且我们有 \(\mathrm{U}_i \cap s_i. \mathrm{U}_i \subset \mathrm{V}_i \cap \mathrm{V}_i' = \emptyset\)。令 \(\mathrm{U}' = \mathrm{U}_0 \cap \mathrm{U}_1 \cap \ldots \cap \mathrm{U}_n\)；\(\mathrm{U}'\) 是开的，包含 \(x\)，并且 \(\mathrm{U}' \cap s. \mathrm{U}' = \emptyset\) 对 \(s \notin \mathbf{K}_x\) 成立。令 \(\mathrm{U} = \bigcap_{t \in \mathbf{K}_x} t. \mathrm{U}'\)，我们最终得到一个开集，它在 \(\mathbf{K}_x\) 下稳定，包含 \(x\)，并且 \(\mathrm{U} \cap s. \mathrm{U} = \emptyset\) 对 \(s \notin \mathbf{K}_x\) 成立：\(\mathrm{U}\) 就是所要求的开集。

典范映射 \(\mathbf{U} / \mathbf{K}_x\to \mathbf{X} / \mathbf{G}\) 是 \(\mathbf{U} / \mathbf{K}_x\) 到 \(\mathbf{X} / \mathbf{G}\) 中一个开集上的同胚这一事实，由第一章 § 5 第 2 节命题 4 得出，因为 \(\mathbf{U}\) 是开的且由 \(\mathbf{G}\) 定义的等价关系是开的（§ 2，第 4 节，引理 2）。

推论. 若再假设 \(K_{x}=\{e\}\)，则点 x 有一个开邻域 U，使得典范映射 \(X\to X/G\) 在 U 上的限制是 U 到 X/G 的一个开子集上的同胚。

